\stagechapter{H}{Harm tree}

\noindent\textbf{What this stage covers.} How possible harms are derived and organised so that app material can be compared with them: the definition of harm, the typology used, the structure of the tree (harm, capabilities, system components), the origin of each component, and the applicability record per app. The draft thesis builds a reference tree per domain (fitness, fertility) from prior tracker research, with explicit AND/OR conditions following Knabenhans et al., component origins (core functionality, digital context, implementation choice), an applicability record per app and pathway (applicable, not applicable, uncertain, with reason), harm categories taken from Knabenhans et al. with a relationship tag from Citron and Solove, and the statement that the trees are reference models of possible harms, not descriptions of actual systems. The basis of the ``best app'' comparison is decision D12 (open, needs the supervisor).

\section{Definitions: what counts as harm, risk and problem}
\pa{citron2021} define harms as ``injuries, setbacks, losses, or impairments to well-being'' \pc{citron2021}{7}{Harms involve injuries, setbacks, losses, or impairments to well-being}, argue that the requirement of harm has impeded enforcement \pc{citron2021}{1}{The requirement of harm has significantly impeded the enforcement of privacy law}, that small harms aggregate at scale \pc{citron2021}{24}{Privacy harms often involve the aggregation of many small harms to each individual, which is compounded by the aggregation of all these harms to many individuals}, that risk is involved in many harm types \pc{citron2021}{25}{Risk is involved with many different types of privacy harm}, and that harms are contextual \pc{citron2021}{26}{Privacy harms are highly contextual, with the harm depending upon how the data is used, what data is involved, and how the data might be combined with other data}. \pa{knabenhans2026} define harm, after the CyberPeace Institute, as an impairment of an entity's capacity to function as it otherwise would in its usual context, with the present practice as baseline. \pa{brooks2017} (NIST) frames risk as likelihood times impact of a \emph{problematic data action} on individuals, which suits findings expressed as app data actions in a context. \pa{nissenbaum2004} defines a violation as a breach of appropriateness or flow norms of a context and names values (harm prevention, informational inequality, autonomy, freedom, relationships, democracy) that decide whether a breach matters.

\section{Typologies}
\begin{itemize}
\item \emph{Citron and Solove: seven harm types.} ``Our typology groups privacy harms into seven basic types'' \pc{citron2021}{39}{Our typology groups privacy harms into seven basic types}: physical, economic, reputational, psychological, autonomy, discrimination and relationship harms, with severity and aggregation as dimensions \pc{citron2021}{38}{Privacy harms not only differ in type but also in their severity}. Autonomy harms have six subtypes (coercion, manipulation, failure to inform, thwarted expectations, lack of control, chilling effect) \pc{citron2021}{53}{People are either directly denied the freedom to decide or are tricked into thinking that they are freely making choices when they are not}; the ``failure to inform'' and ``thwarted expectations'' subtypes fit the thesis' focus on what policies and interfaces communicate \pc{citron2021}{57}{The harm caused by thwarted expectations involves the undermining of people’s choices, such as breaking promises made about the collection, use, and disclosure of personal data}.
\item \emph{Solove 2006: sixteen harmful activities in four groups} (information collection, processing, dissemination, invasion) as a map of \emph{where in the data flow} a problem arises; \pa{solove2006} is conceptual and legal, useful for organising findings by stage of the data flow.
\item \emph{Threat models.} \pa{deng2011} (LINDDUN) gives seven threat types (linkability, identifiability, non-repudiation, detectability, disclosure of information, content unawareness, policy and consent non-compliance) and threat trees with preconditions; \pa{iwaya2023} mapped every finding on 27 mental-health apps to LINDDUN and found that apps do not explain downstream risks \pc{iwaya2023}{3}{apps do not explain (e.g., in the privacy policy) the risks posed by}, which is the gap the thesis measures.
\item \emph{Functional damage potential.} \pa{dehling2015} rated app clusters on five damage characteristics and derived 12 archetypes from store metadata; \pa{grundy2019} traced data to first, third and ``fourth'' parties and modelled aggregation and re-identification potential.
\end{itemize}

\section{The tree structure and component origin}
\pa{knabenhans2026} build trees from three element kinds: system components, capabilities, harms \pc{knabenhans2026}{6}{We build trees from three kinds of elements: system components, capabilities, and harms}. A tree has three levels: the harm at the top, the capabilities that enable it, and system components at the leaves \pc{knabenhans2026}{6}{A harm tree consists of three levels. At the top of the tree is the harm. On the second level, we have the capabilities that enable this harm. At the leaves, we have system components}. Leaves are combined with AND, OR and ``OR+'' and coloured by origin: inherent to the core functionality, inherent to a digital context, or arising from a concrete implementation, the only class that privacy-enhancing technologies can remove. Their conclusion is that in each tree ``there exists at least one pathway to the harm that does not require red components'' \pc{knabenhans2026}{10}{in each harm tree, there exists at least one pathway to the harm that does not require red components}. They claim no completeness \pc{knabenhans2026}{7}{We do not claim completeness of the harms, capabilities, or system components we identify in our analysis}. \pa{deng2011} supplies the same logic from security engineering: a root threat, OR/AND preconditions and misuse-case narratives.

\section{Domain evidence for the leaves and the harms}
\begin{itemize}
\item \emph{Users' scenarios and concerns.} \pa{gabriele2020} (212 fitness-tracker users, 20 threat scenarios rated for plausibility and likelihood); \pa{cao2024} (183 women, post-\emph{Dobbs}: concern is driven by the recipient of the data, with law enforcement and third parties far above the app company, and by extra data types such as location and intimacy \pc{cao2024}{8}{the party with whom the data is being shared with (Data Sharing) was the most important factor}); \pa{moniz2026} (eight Portuguese participants: emotional strain, targeted content after entering intimate data, commodification, discrimination); \pa{martin2023} (purpose and inference matter more than data type: \pc{martin2023}{64}{how information was used, specifically, whether in service of contextual versus non-contextual ends was even more important for our respondents than what data was collected or with whom the data was shared}).
\item \emph{Inference and ecosystem.} \pa{niksirat2024} (236 studies on wearable trackers: health, mental-state, activity and location inferences; insurers, employers, advertisers and abusive partners as recipients; limited user awareness of third-party access); \pa{grundy2019}, \pa{huckvale2019}, \pa{arkalakis2024} for actors and the gap between declared and observed sharing.
\item \emph{Policy and interface evidence.} \pa{malki2024}, \pa{shipp2020}, \pa{hudig2025} (ten transparency and control mechanisms by lifecycle), \pa{doan2025} (stated risks and benefits, collective harms), \pa{zych2026} (only 6\,\% of assessed mental-health apps met the item ``does the developer make clear risks associated with using the app'' \pc{zych2026}{12}{Does the developer make clear risks associated with using the app}).
\item \emph{Regulatory statements of harm.} \pa{bfarm2023}: users in an illness situation ``must not be exploited''; \pa{wp292018}: controllers should spell out the most important consequences of complex or unexpected processing.
\item \emph{Pathways with a named failure condition (added 23 October 2026).} Endpoint privacy zones in fitness networks hide the start and end of routes, yet the protected place could be inferred for 95.1\,\% of moderately active Strava users (\pa{hassan2018}) and, after the vendors' fixes, still for up to 85\,\% of zones \pc{dhondt2022}{1}{our attack discovers the protected location for up to 85\% of EPZs} (\pa{dhondt2022}). Personality traits can be inferred from steps and heart rate, with advertisers, insurers, employers and governments named as interested parties (\pa{zufferey2023}). For FemTech, \pa{mehrnezhad2024} list intimate partners, family, employers, insurers, advertisers, political and religious organisations and governments as threat actors, and note that the data concern partners and children too; EU law does not name FemTech data.
\end{itemize}

\begin{table}[t]
\centering\small
\caption{Harm frameworks and what each supplies to the tree.}
\label{tab:H-frameworks}
\begin{tabularx}{\textwidth}{@{}P{2.8cm}YY@{}}
\toprule
\textbf{Source} & \textbf{What it supplies} & \textbf{Caution}\\
\midrule
\pa{citron2021} & harm types, severity, aggregation; autonomy subtypes tied to communication & US legal frame; harms listed by courts and cases\\
\pa{knabenhans2026} & tree levels, AND/OR, component origin, harm baseline & digital identity, not health apps; harms derived by the authors\\
\pa{solove2006} & activity stages of the data flow & conceptual\\
\pa{deng2011} & threat and precondition trees, soft-privacy threats & engineering method\\
\pa{brooks2017} & problematic data action, likelihood $\times$ impact & introductory report; no catalogue\\
\pa{nissenbaum2004} & norm breach as the core; six values & no coding procedure\\
\pa{martin2023} & purpose and inference as drivers & vignette survey, US\\
\pa{dehling2015} & damage characteristics from store metadata & English store, 2013 data\\
\pa{hassan2018}, \pa{dhondt2022} & a fitness PET with a quantified failure (location) & attack feasibility, not victims\\
\pa{zufferey2023} & inference capability (personality) and interested parties & one device, volunteers\\
\pa{mehrnezhad2024} & FemTech threat actors; third parties affected & products illustrative, not sampled\\
\bottomrule
\end{tabularx}
\end{table}

\begin{thesisbox}[Implications for this thesis (stage H)]
\begin{enumerate}[leftmargin=1.3em]
\item The draft's structure (harm, capabilities, components; AND/OR; origin; applicability record) is a faithful use of \pa{knabenhans2026}. Record which parts are adapted: domain (health apps instead of digital identity), leaves tagged \emph{observed} versus \emph{inferred} because the thesis sees only statements and interfaces, and harm categories mapped to Citron and Solove types.
\item State the harm definition and baseline once (user's situation without the modelled pathway) and keep ``generic risk'' and ``concrete harm'' separate, as the coding axes do.
\item Source each pathway: the data, adversary, access condition, capability and consequence must each have a cited origin (users' scenarios, inference studies, policy evidence). The domain sources in the draft (\pa{niksirat2024}, \pa{gabriele2020}, \pa{shipp2020}, \pa{malki2024}, \pa{hudig2025}, \pa{cao2024}) match the evidence above; \pa{moniz2026} and \pa{martin2023} can add user-perceived harms and purposes.
\item D12 (basis of the ``best app'' comparison) should be settled by the supervisor; \pa{brooks2017} and \pa{citron2021} both warn that severity is contextual, so a ranking of apps by number of applicable pathways is not a ranking of harm.
\item Use pathways whose protection is known to fail as test cases for Axis C: a privacy-zone claim can be compared with the conditions under which \pa{hassan2018} and \pa{dhondt2022} defeat it, which is what a human-judged ``potential mismatch'' needs.
\end{enumerate}
\end{thesisbox}
